
\section{Perspectives et améliorations futures}

À l’issue de notre projet, plusieurs axes d’amélioration ont été identifiés, tant sur le plan fonctionnel que sur le plan technique. Bien que la version actuelle de notre simulateur remplisse les objectifs fixés, il demeure possible d’enrichir considérablement son fonctionnement pour en faire un outil plus robuste, plus immersif et plus intelligent. Ces perspectives ouvrent la voie à une évolution progressive du système vers un simulateur plus complet, plus interactif et potentiellement réutilisable dans d'autres contextes.

\subsection{Améliorations fonctionnelles}

\begin{itemize}
    \item \textbf{Ajout d’un système de missions ou objectifs :} les intrus pourraient avoir un but (atteindre une zone, voler un objet, sortir de la carte), ce qui rendrait leurs mouvements plus crédibles et plus intéressants à contrer. Cela renforcerait la notion de gameplay.
    
    \item \textbf{Création d’un éditeur de carte :} une interface graphique permettant de créer des cartes personnalisées, de placer obstacles, agents et zones sensibles, offrirait une grande liberté de test aux utilisateurs, et favoriserait l’expérimentation.
    
    \item \textbf{Historique des événements :} intégrer un journal des actions (détection, capture, déplacements) visible dans une zone de log de l’interface. Cela aiderait à l’analyse post-simulation et à la traçabilité des comportements.
    
    \item \textbf{Introduction de rôles spécialisés chez les agents :} par exemple, un gardien fixe (caméra), un éclaireur mobile avec vision plus large, ou un chef qui coordonne les autres agents.
     \item \textbf{Affichage dynamique enrichi :} intégration d’une mini-carte, d’un indicateur de danger, de barres d’état ou d’un zoom adaptatif selon les interactions.

    \item \textbf{Paramétrage complet au lancement :} possibilité de choisir tous les paramètres depuis une fenêtre de configuration (taille de la grille, nombre d’agents, fréquence de rafraîchissement, vitesse de jeu).

    \item \textbf{Système de relecture (replay) :} enregistrement des décisions et positions des agents à chaque tour pour permettre une relecture, un partage ou une analyse différée de la simulation.

    \item \textbf{Gestion avancée des collisions :} prise en charge des blocages potentiels entre agents, avec une logique d’attente, de priorité ou de contournement.

    \item \textbf{Mode spectateur et caméras libres :} possibilité pour l’utilisateur de se déplacer librement dans la grille ou de suivre un agent en particulier de manière fluide.
\end{itemize}


\subsection{Améliorations techniques}

\begin{itemize}
    \item \textbf{Affichage dynamique enrichi :} intégration d’une mini-carte, d’un indicateur de danger, de barres d’état ou d’un zoom adaptatif selon les interactions.

    \item \textbf{Paramétrage complet au lancement :} possibilité de choisir tous les paramètres depuis une fenêtre de configuration (taille de la grille, nombre d’agents, fréquence de rafraîchissement, vitesse de jeu).

    \item \textbf{Système de relecture (replay) :} enregistrement des décisions et positions des agents à chaque tour pour permettre une relecture, un partage ou une analyse différée de la simulation.

    \item \textbf{Gestion avancée des collisions :} prise en charge des blocages potentiels entre agents, avec une logique d’attente, de priorité ou de contournement.

    \item \textbf{Mode spectateur et caméras libres :} possibilité pour l’utilisateur de se déplacer librement dans la grille ou de suivre un agent en particulier de manière fluide.
\end{itemize}


\subsection{Compétences techniques consolidées}

Tout au long de ce projet, nous avons mobilisé et renforcé plusieurs savoir-faire essentiels en informatique, notamment :

\begin{itemize}
    \item la modélisation orientée objet (classes, héritage, encapsulation) ;
    \item la conception logicielle modulaire avec le paradigme MVC ;
    \item la programmation concurrente avec les threads et la synchronisation ;
    \item la manipulation d’interfaces graphiques en Java via Swing ;
    \item la structuration d’un code maintenable et extensible ;
    \item l’utilisation rigoureuse de GitHub pour la gestion de versions en équipe.
\end{itemize}

